% =====================================================================
\section{The system under study}
\label{sec:system}
% Condensed from the conference version (hoang2026eas), rewritten rather than
% reused. Code: backend/app/core/supervision.py, graph.py; agents in
% backend/app/agents/. Do not copy prose from ../main.tex.
% =====================================================================

The system is the review-analysis pipeline and runtime supervisor of the
conference version of this work~\citep{hoang2026eas}. We describe here only what
the study manipulates or measures. Neither the pipeline nor the supervisor was
changed for the study.

\subsection{A staged pipeline}
\label{sec:system-pipeline}

The pipeline analyses the reviews of one business in four stages. Each stage is
an LLM agent whose output must validate against a typed contract.
\emph{Analysis} reads the raw reviews and labels each one with an overall
sentiment and a set of aspects drawn from a fixed seven-category taxonomy,
each with its own polarity. \emph{Reasoning} groups these labels into
\emph{patterns}, each consisting of an aspect, a description, a frequency and
a non-empty list of review identifiers cited as evidence. It also proposes
\emph{root causes} that link a pattern to an explanation. \emph{Strategy} turns
patterns and root causes into prioritised recommendations, each consisting of
an issue, an action and an expected impact. \emph{Report} composes the
document a user reads. It restates root causes and recommendations, together
with the business's name and the sample size, which the pipeline supplies as
trusted values.

The pipeline is \emph{closed}. Apart from analysis, which alone reads the raw
reviews, every stage reads only outputs that earlier stages produced in the
same run. No stage has a scratchpad, a retrieval call, a tool or a memory of
earlier runs. Closure gives every claim a checkable provenance. A claim in a
stage's output either derives from upstream outputs or was introduced by that
stage, because information has no other way in. Groundedness is therefore
decidable by consistency checking against upstream state, with no labels and no
second model. We call this the \emph{derivability} property. It licenses
detecting the absence of provenance, not falsity: an output can be faithfully
derived from an upstream error.

\subsection{The supervisor: measure, then decide}
\label{sec:system-supervisor}

Control passes through a single orchestrator node after every stage. The
orchestrator calls two deterministic functions. \code{measure()} computes facts
about the stage's output relative to upstream state and applies no thresholds.
\code{decide()} maps those facts to one of four verdicts by fixed rules:
\verdict{proceed}, \verdict{proceed\_with\_warning}, \verdict{retry} or
\verdict{halt}. Table~\ref{tab:checks} lists the checks. The supervisor makes no
model call. Its verdicts are reproducible, cost nothing at the margin, and can
be recomputed offline from stored facts, which is what makes the replays and
re-scorings in this study possible.

\begin{table}[tbp]
\centering
\caption{The supervisor's checks at the three stages the study injects into.
Each check resolves against upstream state only. An analysis-stage check
(contradiction between a review's star rating and its predicted sentiment)
exists but is not exercised here.}
\label{tab:checks}
\footnotesize
\begin{tabularx}{\linewidth}{@{}l Y L{2.9cm} l@{}}
\toprule
Stage & Measured fact & Check class & Verdict \\
\midrule
Reasoning & a cited evidence identifier does not resolve to any analysed
  review & unresolvable evidence id & \verdict{retry} \\
 & a cited review does not mention the pattern's aspect & misaligned aspect &
  \verdict{retry} \\
 & the claimed frequency differs from the recomputed one by more than 0.05 &
  frequency divergence & warning \\
Strategy & a recommendation's issue has character-level similarity below
  0.35 to every upstream pattern and root cause & untraceable recommendation &
  \verdict{retry} \\
Report & a root cause or recommendation is absent from the upstream outputs &
  report containment & \verdict{retry} \\
 & the business name or sample size differs from the trusted values &
  trusted metadata & \verdict{retry} \\
\bottomrule
\end{tabularx}
\end{table}

A retry re-runs the stage. The retry cap is applied to the verdict, not inside
the checks. Once a stage has been retried twice, a further \verdict{retry}
becomes \verdict{halt} at the analysis, reasoning and report stages, and
\verdict{proceed\_with\_warning} at strategy, which is the one non-critical
stage. Each stage therefore runs at most three times, and the run's length is
bounded by the routing function rather than by the agents eventually
succeeding.

\subsection{Retry feedback}
\label{sec:system-feedback}

On a retry, the measured fact that triggered the verdict is written into the
re-run stage's prompt. The production template, which is arm \arm{A1} of this
study, reads in full:

\begin{prompttext}
A deterministic quality check rejected your previous answer to this task:\\[2pt]
\{feedback\}\\[2pt]
Produce a corrected answer that fixes exactly these issues. Return only valid
JSON.
\end{prompttext}

\noindent Here \code{\{feedback\}} is the supervisor's own measurement. For an
\op{O1} defect from the study it read: \emph{``your previous output cited
unverifiable evidence. pattern `A recurring concern among customers is that the
prices are t': evidence ids not found in the analysis output:
[`6JsbAclWpameErN3DEiFA']''}. That message names the check, the offending item
and the measured value. It does not say what the correct value is. The
production retry appends this message to the stage's original prompt. It does
not include the previous output, so production retries at placement \arm{P0}
in the terms of Section~\ref{sec:design-arms}. Schema validation failures use
the same channel with the validator's error message. The two layers share one
mechanism and differ only in their payload.

\subsection{Recompute-and-overwrite}
\label{sec:system-overwrite}

One fact is corrected rather than retried. A pattern's frequency is the share of
analysed reviews that express its aspect negatively. The model is asked to
state it, but the supervisor can compute it exactly from the analysis output.
The supervisor therefore overwrites the claimed value on every accepted
reasoning output. It warns when the divergence exceeded 0.05, so the
discrepancy stays on record, and never retries. The conference version stated
the principle behind this design (a quantity the checker can compute should not
be asked of the model) but had no counterfactual for it. RQ5 supplies one.

\subsection{Why the pipeline is also a fault-injection testbed}
\label{sec:system-testbed}

The property that makes verification free also makes the pipeline replayable.
Because each stage's input consists entirely of earlier outputs, a stored run
fixes every stage's input exactly. A stage can be re-executed in isolation from
the stored state, a corrupted output can be substituted for a clean one, and
the supervisor can re-measure the result without re-running anything upstream.
The study uses this in three ways. It injects defects into stored clean
outputs; it replays real defects, after first checking that each replay
reproduces the recorded retry exactly; and it re-runs only the stages
downstream of a repaired output.
